\begin{abstract}
The Internet of Vehicles (IoV) is a core part of the Internet of Things (IoT) that connects modern vehicles, roadside units, and traffic systems.
Vehicles need wireless communication to exchange safety alerts and road conditions in real time.
However, connected vehicles face three major practical challenges: selecting the right radio communication technology, sharing sensor data for cooperative perception, and protecting driver privacy while keeping communication secure.
This survey provides a clear, sequential analysis of these three challenges.
First, we compare Cellular V2X (C-V2X: LTE-V2X, 5G NR-V2X) and IEEE 802.11 standards (802.11p, 802.11bd) in terms of medium access, latency, throughput, and global commercial deployment status.
Second, we examine cooperative perception, where vehicles share camera, LiDAR, and radar observations to see beyond line-of-sight obstacles, and we review the strict requirements for low latency and high reliability.
Third, we analyze Vehicular Public Key Infrastructure (V-PKI) and pseudonym certificate systems, evaluating how multi-authority trust, butterfly key expansion, and revocation methods protect privacy while preventing Sybil and tracking attacks.
Finally, we present open challenges in congestion control, edge computing, and misbehavior detection, and outline the roadmap toward 6G-enabled vehicular networks.
\end{abstract}

\begin{IEEEkeywords}
Internet of Vehicles (IoV), Vehicle-to-Everything (V2X), C-V2X, IEEE 802.11bd, Cooperative Perception, Pseudonym Certificates, V-PKI, Intelligent Transportation Systems.
\end{IEEEkeywords}
